\documentclass[UTF8,zihao=-4]{ctexart}

\usepackage[a4paper,margin=2.1cm]{geometry}
\usepackage{xcolor}
\usepackage{enumitem}
\usepackage{hyperref}
\usepackage{parskip}

\definecolor{TitleBlue}{HTML}{1F4E79}
\definecolor{SoftGray}{HTML}{F4F6F8}
\definecolor{TextGray}{HTML}{333333}

\hypersetup{
  colorlinks=true,
  linkcolor=TitleBlue,
  urlcolor=TitleBlue
}

\setlist[itemize]{leftmargin=2em,itemsep=0.2em,topsep=0.2em}
\setlength{\parindent}{0pt}
\renewcommand{\baselinestretch}{1.12}

\newcommand{\principle}[1]{%
  \section*{\color{TitleBlue}#1}
}

\newcommand{\labeltext}[1]{%
  \textbf{#1}
}

\newenvironment{answerbox}{%
  \begin{quote}\small\color{TextGray}
}{%
  \end{quote}
}

\title{\textbf{设计原则速记}}
\author{}
\date{}

\begin{document}
\maketitle

\section*{\color{TitleBlue}总口诀}

\begin{answerbox}
内部暴露 -> 信息隐藏、迪米特法则

职责混乱 -> SRP、高内聚差

直接依赖实现 -> DIP

新增规则总改老代码 -> OCP

接口太大 -> ISP

替换实现出错 -> LSP

包互相依赖 -> ADP

重复逻辑 -> DRY

设计太复杂 -> KISS

提前做没用功能 -> YAGNI
\end{answerbox}

\principle{信息隐藏}

\labeltext{核心：}
模块只暴露必要接口，隐藏内部数据结构和实现细节。

\labeltext{题干信号：}
\begin{itemize}
  \item 直接访问内部对象
  \item 暴露内部数据结构
  \item 字段变化影响多个包
\end{itemize}

\labeltext{答题句：}
\begin{answerbox}
违反信息隐藏原则，因为内部数据结构和变化点没有被封装，其他模块可以直接依赖内部实现。
\end{answerbox}

\principle{SRP：单一职责原则}

\labeltext{核心：}
一个类或包只负责一类职责，只有一个主要变化理由。

\labeltext{题干信号：}
\begin{itemize}
  \item 职责边界不清
  \item 一个类或包做太多事
  \item 业务规则、通知、资源状态、用户权限混在一起
\end{itemize}

\labeltext{答题句：}
\begin{answerbox}
违反 SRP，因为包和类承担多个职责，存在多个变化理由。
\end{answerbox}

\principle{OCP：开闭原则}

\labeltext{核心：}
对扩展开放，对修改关闭。

\labeltext{题干信号：}
\begin{itemize}
  \item 新增规则要修改原有主流程
  \item 变化点没有封装
  \item 每加一种功能都要改旧代码
\end{itemize}

\labeltext{答题句：}
\begin{answerbox}
违反 OCP，因为新增业务规则时需要修改已有主流程，而不是通过新增扩展类完成。
\end{answerbox}

\labeltext{例子：}
\begin{answerbox}
预约冲突检测规则可以抽成 ConflictCheckStrategy。新增时间段冲突、容量冲突、权限冲突、设备状态冲突时，新增策略实现类。
\end{answerbox}

\principle{LSP：里氏替换原则}

\labeltext{核心：}
子类对象应该可以替换父类对象，并且程序行为不被破坏。

\labeltext{题干信号：}
\begin{itemize}
  \item 子类替换父类后行为不一致
  \item 实现类不遵守接口语义
  \item 换一个实现后业务出错
\end{itemize}

\labeltext{答题句：}
\begin{answerbox}
违反 LSP，因为子类或实现类不能在不破坏原有行为的情况下替换父类或接口。
\end{answerbox}

\labeltext{例子：}
\begin{answerbox}
ReservationRepository 的内存实现和数据库实现都应该遵守 save 的语义。

按资源和时间段查询已有预约的方法，也应该在不同实现中保持相同语义。
\end{answerbox}

\principle{ISP：接口隔离原则}

\labeltext{核心：}
接口要小而专一，不要让调用方依赖自己不用的方法。

\labeltext{题干信号：}
\begin{itemize}
  \item 接口太大
  \item 一个接口塞很多职责
  \item 调用方被迫依赖不需要的方法
\end{itemize}

\labeltext{答题句：}
\begin{answerbox}
违反 ISP，因为接口职责过大，调用方依赖了自己不需要的方法。
\end{answerbox}

\labeltext{例子：}
\begin{answerbox}
不要把预约、工单、通知、用户、资源都塞进 CampusService。应该拆成 ReservationService、WorkOrderService、NotificationGateway、UserService、ResourceQueryService。
\end{answerbox}

\principle{DIP：依赖倒置原则}

\labeltext{核心：}
高层业务不直接依赖低层实现，双方都依赖抽象接口。

\labeltext{题干信号：}
\begin{itemize}
  \item 高层业务直接依赖具体实现
  \item 直接 new 低层对象
  \item 直接访问低层内部对象
\end{itemize}

\labeltext{答题句：}
\begin{answerbox}
违反 DIP，因为高层业务直接依赖低层实现对象，而不是依赖稳定接口。
\end{answerbox}

\labeltext{正确方向：}
\begin{answerbox}
ReservationService -> ReservationRepository 接口

MysqlReservationRepository -> 实现 ReservationRepository
\end{answerbox}

\labeltext{错误方向：}
\begin{answerbox}
ReservationService -> MysqlReservationRepository
\end{answerbox}

\principle{Law of Demeter：迪米特法则}

\labeltext{核心：}
一个对象只和直接相关的对象通信，不要穿透别人内部结构。

\labeltext{题干信号：}
\begin{itemize}
  \item 跨对象访问内部结构
  \item 链式访问过多
  \item 知道太多别的对象细节
\end{itemize}

\labeltext{答题句：}
\begin{answerbox}
违反 Law of Demeter，因为对象跨越边界访问其他对象的内部结构，知道了过多实现细节。
\end{answerbox}

\labeltext{坏味道：}
\begin{verbatim}
a.getB().getC().getD()
\end{verbatim}

\principle{ADP：无环依赖原则}

\labeltext{核心：}
包之间的依赖关系不能形成环。

\labeltext{题干信号：}
\begin{itemize}
  \item 包互相依赖
  \item 循环依赖
  \item A 依赖 B，B 又依赖 A
\end{itemize}

\labeltext{答题句：}
\begin{answerbox}
违反 ADP，因为包依赖关系中存在环，修改一个包会牵连其他包。
\end{answerbox}

\labeltext{错误例子：}
\begin{answerbox}
reservation -> resource

resource -> reservation
\end{answerbox}

\principle{DRY：不要重复}

\labeltext{核心：}
同一知识或规则只保留一份。

\labeltext{题干信号：}
\begin{itemize}
  \item 同一逻辑多处复制
  \item 规则散落在多个地方
  \item 改一处还要同步改多处
\end{itemize}

\labeltext{答题句：}
\begin{answerbox}
违反 DRY，因为相同规则在多个位置重复出现，导致修改成本和不一致风险增加。
\end{answerbox}

\principle{KISS：保持简单}

\labeltext{核心：}
能简单解决，就不要设计得过度复杂。

\labeltext{题干信号：}
\begin{itemize}
  \item 设计过度
  \item 简单问题复杂化
  \item 结构很绕但收益不明显
\end{itemize}

\labeltext{答题句：}
\begin{answerbox}
违反 KISS，因为设计复杂度超过当前问题需要，增加了理解和维护成本。
\end{answerbox}

\principle{YAGNI：不要提前设计用不到的东西}

\labeltext{核心：}
当前没有明确需求的功能，不要提前实现。

\labeltext{题干信号：}
\begin{itemize}
  \item 为了未来功能提前写大量代码
  \item 当前需求用不到
  \item 过早抽象
\end{itemize}

\labeltext{答题句：}
\begin{answerbox}
违反 YAGNI，因为实现了当前需求并不需要的功能或抽象。
\end{answerbox}

\principle{高内聚、低耦合}

\labeltext{核心：}
高内聚是一个模块内部的内容围绕同一职责组织。低耦合是模块之间尽量少知道彼此细节。

\labeltext{题干信号：}
\begin{itemize}
  \item 包之间互相直接访问内部对象 -> 低耦合被破坏
  \item 多个职责混在不同包中 -> 高内聚被破坏
  \item 共享可变全局配置 -> 公共耦合
\end{itemize}

\labeltext{第六题可用答案：}
\begin{answerbox}
从耦合角度看，各包互相直接访问内部对象，导致高访问耦合和循环依赖；多个包共享同一个可变全局配置对象，形成公共耦合。

从内聚角度看，预约、资源、工单、通知和用户权限等职责混在不同包中，职责边界不清，导致模块内聚性下降。
\end{answerbox}

\end{document}
